\chapter{Anforderungen modellieren}

\begin{lernziele}
Nach diesem Kapitel können Sie Anforderungen verständlich, eindeutig und prüfbar formulieren. Sie können Anforderungen nummerieren, strukturieren, in funktionale und nicht-funktionale Anforderungen unterscheiden und mit SysML-Beziehungen verbinden. Außerdem verstehen Sie, wie Anforderungen im Roboterprojekt von Stakeholderbedürfnissen bis zu Testfällen nachverfolgt werden.
\end{lernziele}

\section{Einleitung}
Anforderungen sind eine der wichtigsten Grundlagen eines Systemmodells. Sie beschreiben, was das System leisten muss und unter welchen Bedingungen es dies tun soll. Ohne gute Anforderungen ist unklar, welche Architektur sinnvoll ist und woran der Projekterfolg gemessen wird.

Beim autonomen Indoor-Roboter klingen viele Wünsche zunächst einfach: Der Roboter soll sicher fahren, Hindernisse erkennen, Personen erkennen, lange genug fahren und leicht bedienbar sein. Für die Entwicklung reicht das nicht. Solche Wünsche müssen in Anforderungen übersetzt werden, die eindeutig und später überprüfbar sind.

\section{Was ist eine gute Anforderung?}
Eine gute Anforderung ist:

\begin{itemize}
    \item verständlich
    \item eindeutig
    \item notwendig
    \item realistisch
    \item prüfbar
    \item möglichst lösungsneutral
\end{itemize}

\enquote{Der Roboter soll gut navigieren} ist keine gute Anforderung. Was bedeutet \enquote{gut}? Schnell? Genau? Sicher? Ohne Kollision? In unbekannten Räumen? Eine bessere Formulierung lautet:

\begin{quote}
Der Roboter muss in einem bekannten Innenraum ein ausgewähltes Ziel autonom anfahren und dabei statische Hindernisse im Fahrweg vermeiden.
\end{quote}

Auch diese Anforderung kann später noch präziser werden, aber sie ist deutlich verständlicher.

\section{Vom Bedürfnis zur Anforderung}
Stakeholder äußern häufig Bedürfnisse. Diese Bedürfnisse sind wichtig, aber noch keine technischen Anforderungen.

\begin{center}
\begin{tabular}{p{0.38\textwidth}p{0.50\textwidth}}
\toprule
\textbf{Bedürfnis} & \textbf{Mögliche Anforderung} \\
\midrule
Der Roboter soll niemanden gefährden. & Der Roboter muss bei einem Hindernis im Abstand von 0,5 m seine Geschwindigkeit reduzieren und bei 0,2 m stoppen. \\
Der Roboter soll einfach bedienbar sein. & Ein Transportauftrag muss mit höchstens drei Eingaben gestartet werden können. \\
Der Roboter soll nicht liegen bleiben. & Der Roboter muss bei niedrigem Akkustand automatisch zur Ladestation zurückkehren. \\
Wartung soll einfach sein. & Der Roboter muss Sensorstatus, Akkustatus und Fehlermeldungen anzeigen. \\
\bottomrule
\end{tabular}
\end{center}

\section{Beispielanforderungen}
\begin{longtable}{p{0.16\textwidth}p{0.69\textwidth}}
\toprule
ID & Anforderung\\
\midrule
REQ-001 & Der Roboter muss autonom in einem bekannten Innenraumbereich navigieren können.\\
REQ-002 & Der Roboter muss Hindernisse im Fahrbereich erkennen.\\
REQ-003 & Der Roboter muss erkannte Hindernisse umfahren oder rechtzeitig stoppen.\\
REQ-004 & Der Roboter muss Personen im Kamerabild erkennen können.\\
REQ-005 & Der Roboter muss seinen Akkustand überwachen.\\
REQ-006 & Der Roboter muss bei niedrigem Akkustand selbstständig zur Ladestation fahren.\\
REQ-007 & Der Roboter muss seinen Betriebszustand melden.\\
REQ-008 & Der Roboter muss einen sicheren Zustand einnehmen, wenn ein kritischer Sensor ausfällt.\\
REQ-009 & Der Roboter muss Transportaufträge über eine einfache Benutzerschnittstelle entgegennehmen.\\
REQ-010 & Der Roboter muss relevante Fehler für Wartungspersonal nachvollziehbar protokollieren.\\
\bottomrule
\end{longtable}

\section{Funktionale und nicht-funktionale Anforderungen}
Funktionale Anforderungen beschreiben, was das System tun soll. Beispiele sind Navigation, Hinderniserkennung, Akkustandsüberwachung oder Statusmeldung.

Nicht-funktionale Anforderungen beschreiben Eigenschaften des Systems. Beispiele sind Geschwindigkeit, Genauigkeit, Verfügbarkeit, Sicherheit, Wartbarkeit, Datenschutz oder Energieverbrauch.

Beide Arten sind wichtig. Ein Roboter, der alle Funktionen besitzt, aber zu langsam, unsicher oder unwartbar ist, erfüllt das Projektziel nicht.

\section{Anforderungen strukturieren}
In größeren Projekten werden Anforderungen hierarchisch strukturiert. Eine oberste Anforderung kann in mehrere konkrete Anforderungen zerlegt werden.

Beispiel:

\begin{itemize}
    \item REQ-SAFE-001: Der Roboter muss sicher in der Nähe von Menschen betrieben werden.
    \item REQ-SAFE-002: Der Roboter muss Hindernisse im Fahrweg erkennen.
    \item REQ-SAFE-003: Der Roboter muss bei Gefahr stoppen.
    \item REQ-SAFE-004: Der Roboter muss bei Sensorausfall einen sicheren Zustand einnehmen.
\end{itemize}

Die allgemeine Sicherheitsanforderung wird dadurch handhabbarer.

\section{SysML-Beziehungen}
Wichtige Beziehungen in SysML sind:

\begin{itemize}
    \item \texttt{derive}: Eine Anforderung wird aus einer anderen abgeleitet.
    \item \texttt{satisfy}: Ein Modellelement erfüllt eine Anforderung.
    \item \texttt{verify}: Ein Testfall überprüft eine Anforderung.
    \item \texttt{trace}: Eine allgemeine Nachverfolgungsbeziehung.
    \item \texttt{refine}: Ein Element beschreibt eine Anforderung genauer.
\end{itemize}

\begin{praxisbox}
REQ-002 \enquote{Hindernisse erkennen} wird durch das \texttt{PerceptionSystem} erfüllt. Der Testfall \texttt{TC-002 Hinderniserkennung} verifiziert diese Anforderung. Eine \texttt{trace}-Beziehung kann zusätzlich zeigen, dass der ROS2-Node \texttt{obstacle\_detector\_node} zur Umsetzung gehört.
\end{praxisbox}

Abbildung~\ref{fig:kap06-traceability-req002} stellt diese Traceability-Kette grafisch dar: \texttt{satisfy} und \texttt{verify} zeigen als zwei getrennte, eingehende Beziehungen auf das \texttt{PerceptionSystem}, von dort führt \texttt{trace} weiter zum ROS2-Node.

\begin{figure}[H]
  \centering
  \adjustbox{max width=\linewidth,max totalheight=0.6\textheight}{% =====================================================================
% Abbildung: Traceability-Kette am Beispiel REQ-002
% Passend zu Kapitel 6, Abschnitt "SysML-Beziehungen"
%
% Einbindung im Buch, z. B. in chapters/06_anforderungen.tex:
%
%   \begin{figure}[H]
%     \centering
%     % =====================================================================
% Abbildung: Traceability-Kette am Beispiel REQ-002
% Passend zu Kapitel 6, Abschnitt "SysML-Beziehungen"
%
% Einbindung im Buch, z. B. in chapters/06_anforderungen.tex:
%
%   \begin{figure}[H]
%     \centering
%     % =====================================================================
% Abbildung: Traceability-Kette am Beispiel REQ-002
% Passend zu Kapitel 6, Abschnitt "SysML-Beziehungen"
%
% Einbindung im Buch, z. B. in chapters/06_anforderungen.tex:
%
%   \begin{figure}[H]
%     \centering
%     \input{figures/fig_kap06_traceability_req002}
%     \caption{Traceability-Kette von REQ-002 ueber \texttt{satisfy} und
%       \texttt{verify} zum \texttt{PerceptionSystem} und per \texttt{trace}
%       weiter zum ROS2-Node.}
%     \label{fig:kap06-traceability-req002}
%   \end{figure}
%
% Benoetigte Pakete (im Buch bereits in main.tex geladen): tikz
% Benoetigte TikZ-Bibliotheken: positioning, arrows.meta, shadows, calc
% =====================================================================

\input{figures/fig_kap06_traceability_req002.tikz}

%     \caption{Traceability-Kette von REQ-002 ueber \texttt{satisfy} und
%       \texttt{verify} zum \texttt{PerceptionSystem} und per \texttt{trace}
%       weiter zum ROS2-Node.}
%     \label{fig:kap06-traceability-req002}
%   \end{figure}
%
% Benoetigte Pakete (im Buch bereits in main.tex geladen): tikz
% Benoetigte TikZ-Bibliotheken: positioning, arrows.meta, shadows, calc
% =====================================================================

\begin{tikzpicture}[
    every node/.style={font=\small},
    req/.style={
        rectangle, rounded corners=1.5pt, draw=black!70, thick,
        minimum width=34mm, minimum height=12mm,
        align=center, fill=blue!12, drop shadow
    },
    part/.style={
        rectangle, rounded corners=1.5pt, draw=green!40!black, thick,
        minimum width=34mm, minimum height=12mm,
        align=center, fill=green!15, drop shadow
    },
    test/.style={
        rectangle, rounded corners=1.5pt, draw=black!50, thick,
        minimum width=34mm, minimum height=12mm,
        align=center, fill=black!6, drop shadow
    },
    code/.style={
        rectangle, rounded corners=1.5pt, draw=orange!60!black, thick,
        minimum width=40mm, minimum height=12mm,
        align=center, fill=orange!15, drop shadow
    },
    flow/.style={-{Stealth[length=2.2mm]}, thick, black!60},
    lbl/.style={font=\small, fill=white, inner sep=1pt}
]

% --- Knoten ---
\node[req]  (req002)  at (0, 3)     {\texttt{REQ-002}\\Hindernisse erkennen};
\node[part] (percept) at (7, 3)     {\texttt{PerceptionSystem}};
\node[test] (tc002)   at (0, -1.5)  {\texttt{TC-002}\\Hinderniserkennung};
\node[code] (rosnode) at (7, -1.5)  {\texttt{obstacle\_detector\_node}};

% --- Beziehungen ---
\draw[flow] (req002) -- node[lbl, midway]{\texttt{satisfy}} (percept);
\draw[flow] (tc002)  to[bend left=12] node[lbl, near end]{\texttt{verify}} (percept);
\draw[flow] (percept) -- node[lbl, midway]{\texttt{trace}} (rosnode);

% --- Legende ---
\node[req,  minimum width=6mm, minimum height=4mm, below=15mm of tc002, xshift=-8mm, font=\tiny] (leg1) {};
\node[right=1mm of leg1, font=\tiny, align=left] {Anforderung};
\node[part, minimum width=6mm, minimum height=4mm, below=6mm of leg1, font=\tiny] (leg2) {};
\node[right=1mm of leg2, font=\tiny, align=left] {Systemteil (\texttt{part def})};
\node[test, minimum width=6mm, minimum height=4mm, below=6mm of leg2, font=\tiny] (leg3) {};
\node[right=1mm of leg3, font=\tiny, align=left] {Testfall};
\node[code, minimum width=6mm, minimum height=4mm, below=6mm of leg3, font=\tiny] (leg4) {};
\node[right=1mm of leg4, font=\tiny, align=left] {ROS2-Node};

\end{tikzpicture}


%     \caption{Traceability-Kette von REQ-002 ueber \texttt{satisfy} und
%       \texttt{verify} zum \texttt{PerceptionSystem} und per \texttt{trace}
%       weiter zum ROS2-Node.}
%     \label{fig:kap06-traceability-req002}
%   \end{figure}
%
% Benoetigte Pakete (im Buch bereits in main.tex geladen): tikz
% Benoetigte TikZ-Bibliotheken: positioning, arrows.meta, shadows, calc
% =====================================================================

\begin{tikzpicture}[
    every node/.style={font=\small},
    req/.style={
        rectangle, rounded corners=1.5pt, draw=black!70, thick,
        minimum width=34mm, minimum height=12mm,
        align=center, fill=blue!12, drop shadow
    },
    part/.style={
        rectangle, rounded corners=1.5pt, draw=green!40!black, thick,
        minimum width=34mm, minimum height=12mm,
        align=center, fill=green!15, drop shadow
    },
    test/.style={
        rectangle, rounded corners=1.5pt, draw=black!50, thick,
        minimum width=34mm, minimum height=12mm,
        align=center, fill=black!6, drop shadow
    },
    code/.style={
        rectangle, rounded corners=1.5pt, draw=orange!60!black, thick,
        minimum width=40mm, minimum height=12mm,
        align=center, fill=orange!15, drop shadow
    },
    flow/.style={-{Stealth[length=2.2mm]}, thick, black!60},
    lbl/.style={font=\small, fill=white, inner sep=1pt}
]

% --- Knoten ---
\node[req]  (req002)  at (0, 3)     {\texttt{REQ-002}\\Hindernisse erkennen};
\node[part] (percept) at (7, 3)     {\texttt{PerceptionSystem}};
\node[test] (tc002)   at (0, -1.5)  {\texttt{TC-002}\\Hinderniserkennung};
\node[code] (rosnode) at (7, -1.5)  {\texttt{obstacle\_detector\_node}};

% --- Beziehungen ---
\draw[flow] (req002) -- node[lbl, midway]{\texttt{satisfy}} (percept);
\draw[flow] (tc002)  to[bend left=12] node[lbl, near end]{\texttt{verify}} (percept);
\draw[flow] (percept) -- node[lbl, midway]{\texttt{trace}} (rosnode);

% --- Legende ---
\node[req,  minimum width=6mm, minimum height=4mm, below=15mm of tc002, xshift=-8mm, font=\tiny] (leg1) {};
\node[right=1mm of leg1, font=\tiny, align=left] {Anforderung};
\node[part, minimum width=6mm, minimum height=4mm, below=6mm of leg1, font=\tiny] (leg2) {};
\node[right=1mm of leg2, font=\tiny, align=left] {Systemteil (\texttt{part def})};
\node[test, minimum width=6mm, minimum height=4mm, below=6mm of leg2, font=\tiny] (leg3) {};
\node[right=1mm of leg3, font=\tiny, align=left] {Testfall};
\node[code, minimum width=6mm, minimum height=4mm, below=6mm of leg3, font=\tiny] (leg4) {};
\node[right=1mm of leg4, font=\tiny, align=left] {ROS2-Node};

\end{tikzpicture}

}
  \caption{Traceability-Kette von REQ-002 über \texttt{satisfy} und \texttt{verify} zum \texttt{PerceptionSystem} und per \texttt{trace} weiter zum ROS2-Node.}
  \label{fig:kap06-traceability-req002}
\end{figure}

\section{Anforderungen in SysML v2}
In SysML v2 werden Anforderungen – eingebettet in ein \texttt{package} wie bereits aus Kapitel 3 bekannt – textuell mit dem Schlüsselwort \texttt{requirement} definiert:

\begin{lstlisting}
package Requirements {

  requirement REQ_001 {
    doc /* Der Roboter muss in einem bekannten Innenraumbereich
           autonom navigieren koennen. */
  }

  requirement REQ_SAFE {
    doc /* Sicherheitsanforderungen. */
    requirement REQ_SAFE_001 {
      doc /* Der Roboter muss Hindernisse im Fahrweg erkennen. */
      derive REQ_001;
    }
  }
}
\end{lstlisting}

Anders als bei \texttt{part def} und \texttt{part} in Kapitel 8 unterscheidet SysML v2 bei \texttt{requirement} im Grundfall nicht zwischen Typ und Nutzung: Das Schlüsselwort \texttt{requirement} definiert die Anforderung im Beispiel gleichzeitig als solche.

Erstellen Sie in der Datei \texttt{02\_requirements.sysml} erste Requirement-Elemente. Verwenden Sie klare Namen, eine stabile ID und einen prüfbaren Text. Achten Sie darauf, dass Anforderungen nicht nur als Liste existieren, sondern später mit Systemteilen, Verhalten und Tests verbunden werden.

Nutzen Sie Beschreibungsfelder (\texttt{doc /* ... */}), um Annahmen, Begründungen oder offene Fragen zu dokumentieren. Wenn eine Anforderung noch unklar ist, markieren Sie sie bewusst, statt sie stillschweigend als fertig zu behandeln – zum Beispiel mit einem kurzen \texttt{doc /* TBD: ... */}-Vermerk im betroffenen Requirement-Element.

\section{Typische Fehler}
\subsection{Zu allgemeine Anforderungen}
\enquote{Der Roboter soll zuverlässig sein} ist nicht prüfbar. Besser sind konkrete Kriterien.

\subsection{Lösungen als Anforderungen}
\enquote{Der Roboter muss Sensor X verwenden} ist häufig bereits eine Designentscheidung. Besser ist eine lösungsneutrale Fähigkeit, sofern kein konkreter Sensor vorgeschrieben ist.

\subsection{Keine Akzeptanzkriterien}
Eine Anforderung sollte erkennen lassen, wie sie später geprüft wird. Ein Akzeptanzkriterium legt fest, woran sich die Erfüllung konkret zeigt, zum Beispiel: \enquote{Testfall: Bei einem simulierten Hindernis in 0,5\,m Abstand muss die Geschwindigkeit innerhalb 1\,s reduziert werden.} Ohne solche Kriterien entstehen Diskussionen am Projektende.

\subsection{Keine Pflege}
Anforderungen ändern sich. Wichtig ist, Änderungen nachvollziehbar zu halten und betroffene Modellelemente zu prüfen.

\section{Vertiefungsbeispiel: Batteriemanagement mit abgesicherter Rückkehr}
\label{sec:bm-anforderungen}
Die Anforderungen REQ-005 und REQ-006 beschreiben zunächst Akkustandsüberwachung und Rückkehr bei niedrigem Akkustand. Wir verfeinern diese Ausgangslage: Der Roboter soll anhand des bisherigen Verbrauchs abschätzen, ob er den laufenden Transportauftrag noch abschließen und anschließend zur Ladestation zurückkehren kann. Ein fester Akkuschwellwert allein beantwortet diese Frage nicht.

Dieses Beispiel ist ein fachlicher Entwurf. Die folgenden BM-Anforderungen ergänzen die bisherigen Lehrbuchbeispiele; ihre Umsetzung in SysML und ROS2 sowie der Nachweis am Roboter stehen noch aus. Der Entscheidungsablauf folgt in Abschnitt~\ref{sec:bm-ablauf}, die Energiebilanz in Abschnitt~\ref{sec:bm-energie} und die Testzuordnung in Abschnitt~\ref{sec:bm-tests}.

\subsection{Systemgrenze und Annahmen}
Betrachtet werden Akkuzustand, Verbrauchserfassung, Energiebedarfsschätzung, Rückwegabsicherung, Auftragsentscheidung, Navigation und Statusmeldungen. Die Ladestation ist ein externes System. Die Rückkehr endet erst mit bestätigtem Andocken und Ladebeginn.

Als Entwurfsannahme beginnt die Fahrt an einer bekannten Ladestation. Der Roboter zeichnet den gefahrenen Weg und dessen Energieverbrauch auf und kann dem Weg zurück folgen. Diese vorgeschlagenen Fähigkeiten müssen technisch geprüft werden. Ein bereits gefahrener Weg ist weder automatisch in Gegenrichtung befahrbar noch dauerhaft frei. Auch der Rückwegverbrauch darf nicht ungeprüft mit dem Hinwegverbrauch gleichgesetzt werden.

Ein \emph{abgesicherter Rückweg} besitzt einen bekannten Verlauf, ist ausreichend zuverlässig verfolgbar und wurde hinsichtlich Befahrbarkeit und Energiebedarf einschließlich Andocken bewertet. Eine \emph{belastbare Schätzung} verwendet gültige Eingangsdaten und berücksichtigt ihre Unsicherheit. Welche Qualität hierfür genügt, bleibt ein offener Projektparameter.

\begin{hinweisbox}
Sind Restauftrag oder vorgesehener Rückweg nicht abschätzbar, etwa wegen fehlender Kartendaten, darf der Auftrag nicht fortgesetzt werden. Die Rückkehr muss dann durch eine unabhängig abgesicherte Möglichkeit gewährleistet werden, beispielsweise durch eine geeignete Wegaufzeichnung. Eine fehlende Schätzung ist kein geringer Energiebedarf.
\end{hinweisbox}

\subsection{Anforderungen des Vertiefungsbeispiels}
Die IDs BM-01 bis BM-13 bleiben über Anforderungen, Ablauf und Tests hinweg stabil. BM-06 legt bewusst eine Entwurfsentscheidung für dieses Beispiel fest; eine andere technische Rückkehrstrategie wäre gesondert zu bewerten.

\begin{description}[style=nextline]
    \item[BM-01: Startfreigabe] Vor Fahrtbeginn muss der Roboter gültige Akkudaten, Verbrauchserfassung und eine verfügbare Rückkehrstrategie prüfen. Fehlt eine Voraussetzung, darf kein Transportauftrag starten.
    \item[BM-02: Verbrauchsbasierte Schätzung] Der Roboter muss den bisherigen Energieverbrauch erfassen und damit den geschätzten Bedarf aktualisieren. Veränderungen von Last, Fahrprofil oder Umgebung müssen in die Bewertung eingehen.
    \item[BM-03: Entscheidungsgrundlage] Für jede Auftragsentscheidung muss der Roboter nutzbare Restenergie, Bedarf des Restauftrags, anschließenden Rückkehrbedarf und Reserve mit Gültigkeit und Unsicherheit bewerten. Unbekannte Werte dürfen nicht als null gelten.
    \item[BM-04: Auftrag fortsetzen] Der Roboter darf den Auftrag nur fortsetzen, wenn Restauftrag, anschließende Rückkehr und Reserve durch eine belastbare Energieschätzung gedeckt sind. Zusätzlich muss eine Rückkehr vom aktuellen Standort abgesichert bleiben.
    \item[BM-05: Fehlende Abschätzbarkeit] Sind Restauftrag oder anschließender Rückweg nicht abschätzbar, muss der Roboter den Auftrag kontrolliert unterbrechen und eine abgesicherte Rückkehr beginnen. Ein hoher Akkustand hebt diese Regel nicht auf.
    \item[BM-06: Wegaufzeichnung] Der Roboter muss den gefahrenen Weg und zugehörige Verbrauchsdaten für die Rückkehrbewertung aufzeichnen. Lücken und ungültige Abschnitte müssen erkannt und berücksichtigt werden.
    \item[BM-07: Rückwegfreigabe] Ein aufgezeichneter Weg darf nur bei bewerteter Verfolgbarkeit, Befahrbarkeit und ausreichender Energie einschließlich Reserve als abgesicherter Rückweg verwendet werden. Die Aufzeichnung allein genügt nicht.
    \item[BM-08: Laufende Absicherung] Der Roboter muss die Rückkehrabsicherung regelmäßig und bei relevanten Änderungen erneut bewerten. Droht sie verloren zu gehen, muss er die Rückkehr vor Unterschreiten der festgelegten Rückkehrgrenze einleiten. Reaktions- und Unterbrechungszeit sind zu berücksichtigen.
    \item[BM-09: Blockierter Rückweg] Bei einer Blockade muss der Roboter einen alternativen Rückweg bewerten. Er darf ihn nur bei abgesicherter Befahrbarkeit und ausreichender Energie nutzen. Auch Warte- und Planungsenergie müssen berücksichtigt werden.
    \item[BM-10: Keine abgesicherte Rückkehr] Ist kein abgesicherter Rückweg verfügbar, muss der Roboter kontrolliert anhalten und Unterstützung anfordern. Der Auftrag bleibt unterbrochen; die autonome Rückkehr muss als nicht erfüllt gemeldet werden.
    \item[BM-11: Datenausfall] Bei Ausfall von Akkudaten oder Wegaufzeichnung muss der Roboter die verbleibende Rückkehrabsicherung neu bewerten. Nur eine weiterhin belastbar abgesicherte Rückkehr darf fortgesetzt werden; andernfalls gilt BM-10.
    \item[BM-12: Nachvollziehbarer Status] Der Roboter muss Entscheidung, Entscheidungsgrund, Auftragsstatus und Rückkehrstatus nachvollziehbar melden. Fortsetzung, Unterbrechung, Rückkehr und Unterstützungsbedarf müssen unterscheidbar sein.
    \item[BM-13: Rückkehr erfolgreich] Der Roboter darf die Rückkehr erst nach bestätigtem Andocken und Ladebeginn als erfolgreich melden. Das Erreichen der Stationsposition allein genügt nicht.
\end{description}

\subsection{Offene Parameter und Nachweisgrenzen}
Vor einer Implementierung müssen das Verfahren zur Restenergieschätzung, Energiereserve, zulässige Schätzunsicherheit, Rückkehrgrenze, Prüfabstände und Kriterien für gültige Wegdaten festgelegt werden. Ebenso ist zu klären, wie ein kontrolliertes Anhalten bei der vorgesehenen Ladung und Umgebung aussieht. Konkrete Werte werden hier bewusst nicht erfunden.

Damit liegt noch keine technische Rückkehrgarantie vor. Insbesondere die Rückkehr ohne Karte, die Verbrauchsschätzung und unerwartete Wegblockaden benötigen einen Nachweis. Kontrolliertes Anhalten mit Unterstützungsanforderung begrenzt die Folgen eines Fehlers, erfüllt aber nicht die geforderte autonome Rückkehr.

\section{Zusammenfassung}
Anforderungen verbinden Stakeholderbedürfnisse mit Architektur und Tests. Gute Anforderungen sind verständlich, eindeutig und prüfbar. In SysML v2 werden Anforderungen nicht nur gesammelt, sondern über Beziehungen mit Systemteilen, Aktionen, Schnittstellen und Testfällen verbunden.

\section{Kontrollfragen}
\begin{enumerate}
    \item Was macht eine gute Anforderung aus?
    \item Warum ist \enquote{Der Roboter soll gut navigieren} problematisch?
    \item Was ist der Unterschied zwischen funktionalen und nicht-funktionalen Anforderungen?
    \item Wozu dient \texttt{derive}?
    \item Wozu dient \texttt{satisfy}?
    \item Wozu dient \texttt{verify}?
    \item Warum sollten Anforderungen möglichst lösungsneutral sein?
    \item Welche Anforderung könnte aus dem Bedürfnis \enquote{Wartung soll einfach sein} entstehen?
\end{enumerate}

\section{Übungen}
\begin{uebungbox}
\textbf{Übung 1: Anforderungen formulieren}\\
Formulieren Sie fünf zusätzliche Anforderungen: zwei zur Sicherheit, eine zur Bedienung, eine zur Wartung und eine zur Leistung.
\end{uebungbox}

\begin{uebungbox}
\textbf{Übung 2: Anforderungen verbessern}\\
Verbessern Sie die Aussagen \enquote{Der Roboter soll schnell sein}, \enquote{Der Roboter soll zuverlässig sein} und \enquote{Die Bedienung soll intuitiv sein}, sodass daraus prüfbare Anforderungen werden.
\end{uebungbox}

\begin{uebungbox}
\textbf{Übung 3: Beziehungen anlegen}\\
Wählen Sie drei Anforderungen aus der Tabelle. Ordnen Sie jeder Anforderung ein Systemelement (\texttt{part def}) und einen Testfall zu.
\end{uebungbox}

\section{Literaturhinweise}
Anforderungsmodellierung ist ein zentraler Bestandteil von SysML und MBSE. Für den Einstieg eignen sich Delligatti \cite{delligatti2014} sowie Friedenthal, Moore und Steiner \cite{friedenthal2014}. Für reale Projekte sollte zusätzlich Literatur zum Requirements Engineering genutzt werden.
